\chapter{Semantic Schema and Label Space}
This chapter defines the internal data model and the controlled vocabularies used to transform raw, heterogeneous internship listings into a structured semantic knowledge base. We describe the \texttt{Internship} interface as a formal contract for the data pipeline and explain the selection of the ESCO taxonomy as the primary reference for mapping job titles and skills.

\section{Target Schema Definition}
To achieve a unified view of the labor market data, the system requires a rigid target schema that can accommodate the variability of source data while providing the necessary metadata for hybrid retrieval.

\subsection{Structural Requirements}
The schema is designed to support three primary objectives:
\begin{enumerate}
    \item \textbf{Identity and Provenance:} Ensuring unique identification across ingestion cycles and maintaining links to original sources.
    \item \textbf{Faceted Discovery:} Providing standardized labels (categories, locations, skills) for efficient filtering.
    \item \textbf{Semantic Enrichment:} Storing high-dimensional representations (embeddings) and canonical labels to facilitate non-lexical search.
\end{enumerate}

\subsection{The Internship Interface}
The schema is implemented as a TypeScript interface, serving as the blueprint for both the database and the frontend presentation layer.

\section{Taxonomy Selection: The ESCO Framework}
A central goal of this research is to map idiosyncratic titles to a standardized professional taxonomy. We have selected the \textbf{ESCO (European Skills, Competences, Qualifications and Occupations)} framework as our primary label space.

\subsection{Occupations Pillar}
ESCO provides a hierarchical structure for occupations. For the purpose of this system, we focus on the "Level 4" unit groups, which provide a balance between specificity and generalizability. For example, a "React Developer" listing is mapped to the ESCO concept \textit{"software developer"}. This mapping is critical for the title normalization task evaluated in Chapter 11.

\subsection{Skills and Competences Pillar}
The skill lexicon used in the extraction phase is grounded in the ESCO Skills pillar. By using ESCO as a reference, the system avoids the "vocabulary gap" where different employers use synonyms to describe the same competency. The system maps these aliases to a single canonical label, enabling higher recall in skill-based filtering.

\section{Internal Categorization and Mapping}
While ESCO provides the granular depth, the user interface requires a flatter, more navigable set of high-level categories to facilitate efficient browsing.

\subsection{Domain-Specific Categories}
We define a set of top-level categories that reflect the current internship market:
\begin{itemize}
    \item \textbf{Software:} Encompasses web, mobile, and systems engineering.
    \item \textbf{Data/AI:} Includes data science, machine learning, and analytics.
    \item \textbf{Design:} Covers UI/UX, graphic design, and industrial design.
    \item \textbf{Healthcare:} Specialized roles in pharmacy, medicine, and health tech.
    \item \textbf{Business:} A broad category for marketing, management, and finance.
\end{itemize}

\subsection{Heuristic Mapping Logic}
To bridge the gap between raw titles and these categories, the system utilizes a keyword-based anchor strategy. This serves as the first layer of the semantic pipeline. Job titles are analyzed for specific keywords:
\begin{itemize}
    \item \textbf{Software:} "developer," "engineer," or "software."
    \item \textbf{Data/AI:} "data," "analyst," "ai," or "scientist."
    \item \textbf{Design:} "design," "ux," or "ui."
\end{itemize}
This lightweight approach ensures that thousands of listings can be categorized in milliseconds without the computational overhead of deep learning inference for simple classification tasks.

\section{Addressing the Semantic Gap}
The chosen label space specifically addresses the "semantic gap" by moving from lexical tokens to conceptual entities. By defining the schema around canonical identifiers rather than raw strings, the system prepares the data for the hybrid search infrastructure discussed in Chapter 9, where queries are matched not just on text, but on their proximity within the defined semantic space.
